\textbf{Volatile storage.} The queue resides entirely in SRAM, so a power interruption discards every
waiting record and the chamber falls back to the manual process the system replaced. Where grid supply
is intermittent this is the most serious functional limitation of the current build. The ATmega32
provides 1~KB of internal EEPROM, enough to mirror a working queue of the depth SRAM supports, and
writing each record through on enqueue would make the queue survive a reset with no added components.
An SD card would additionally permit a retained visit history, at the cost of an SPI interface and a
filesystem layer.

\textbf{Strict FCFS ordering.} The queue implements first-come, first-served with no override. This
is the fair default and the behaviour the system was built to enforce, but clinical practice
sometimes requires otherwise, whether for an acutely unwell patient or on the accessibility grounds
motivating the priority scheme of Chin et al.~\cite{arpn2019}. A second console control promoting a
selected record would address this, and any such mechanism must stay visible to waiting patients to
avoid undermining the transparency that motivated automating the queue.

\textbf{Corners tethered by cable.} Both corners wire back to one controller, so their separation is
bounded by cable length and they cannot occupy different rooms as a larger chamber might require.
The reserved UART lines and the 38400~bps operating point established in
Section~\ref{sec:calibration} exist to lift this restriction, since a Bluetooth or Wi-Fi module on
that port would let the console sit in a separate room and would support mobile waiting-list
notification. Any such extension must preserve the property that the offline path keeps working
when the link is down.

\textbf{Accuracy and scale.} Section~\ref{sec:results} identifies three inexpensive corrections, namely
setting the ADC prescaler to bring the conversion clock into the datasheet's 50--200~kHz window,
storing a per-unit calibration offset for the LM35, and deriving the pulse sampling interval from a
hardware timer rather than delay loops. None requires a different microcontroller or sensor. The
architecture also extends to a multi-room variant with one booth feeding several consoles, where the
queue structure generalises directly but the pin budget does not. Placing all consoles on the shared
I\textsuperscript{2}C bus absorbs that growth. Finally, as Section~\ref{sec:threats} notes, guided
self-registration has not been tested with patients, and a deployment study would test the claim this
design rests on most heavily.
